\chapter{Gonorrhea Testing}

\section*{What Is This Test?}
Gonorrhea testing detects \textit{Neisseria gonorrhoeae}, a gram-negative intracellular diplococcus and the causative agent of gonorrhea, the second most commonly reported bacterial sexually transmitted infection in the United States with approximately 700,000 reported cases annually (likely an underestimate). \textit{N. gonorrhoeae} infects the columnar and transitional epithelium of the urogenital tract (urethra, endocervix), rectum, pharynx, and conjunctiva. Clinical manifestations include urethritis (dysuria, purulent discharge) in men; cervicitis (often asymptomatic, but may present with vaginal discharge, intermenstrual bleeding, dyspareunia) in women; pelvic inflammatory disease (PID) with risk of tubal factor infertility and ectopic pregnancy; pharyngitis (usually asymptomatic); proctitis; and conjunctivitis (especially neonatal ophthalmia neonatorum). Disseminated gonococcal infection (DGI) presents with the classic triad of tenosynovitis, dermatitis (pustular or petechial lesions on extremities), and polyarthralgia/septic arthritis (typically involving wrists, knees, ankles). The CDC recommends NAAT for detection of \textit{N. gonorrhoeae} from genital (vaginal swab, endocervical swab, first-catch urine) and extragenital (rectal, pharyngeal) sites. Antimicrobial resistance in \textit{N. gonorrhoeae} is a critical public health threat with resistance reported to all classes of antibiotics historically used for treatment; ceftriaxone remains the only CDC-recommended first-line therapy.

\section*{Alternative Names}
GC Test; Gonorrhea NAAT; Gonorrhea PCR; N. gonorrhoeae Culture; GC Culture and Sensitivity; Gonococcus Test; Gonorrhea DNA Probe

\section*{Principle}
Nucleic acid amplification tests (NAATs) are the gold standard: Real-time PCR (Roche cobas CT/NG, Abbott RealTime CT/NG, Hologic Aptima Combo 2, Cepheid GeneXpert CT/NG) detecting \textit{N. gonorrhoeae}-specific DNA or ribosomal RNA targets (porA pseudogene, Opa genes, 16S rRNA). Sensitivity 95--99\% for urogenital specimens, 90--95\% for pharyngeal and rectal specimens. Specificity >99\%. NAATs are FDA-cleared for endocervical, vaginal, urine, and in some assays, pharyngeal and rectal specimens. Results available in 1--4 hours. Culture: Specimens inoculated on modified Thayer-Martin (MTM) or New York City (NYC) selective agar containing vancomycin, colistin, nystatin, and trimethoprim to suppress normal flora. Incubated at 35--37\textdegree{}C in 5--7\% CO$_2$ for 48--72 hours. Colonies appear as small, gray, oxidase-positive gram-negative diplococci. Identification by carbohydrate utilization (glucose positive; maltose, lactose, sucrose negative), MALDI-TOF, or DNA probe. Culture has sensitivity of 85--95\% for symptomatic urethral infection but lower for asymptomatic, pharyngeal, and rectal infections. The critical advantage of culture is antimicrobial susceptibility testing (AST). Gram stain: Urethral discharge Gram stain showing gram-negative intracellular diplococci within PMNs is highly sensitive (>95\%) and specific (>99\%) in symptomatic men. In endocervical specimens, sensitivity is only 50--70\% and Gram stain is inadequate for screening. Gram stain of pharyngeal and rectal specimens should not be used.

\section*{History}
The organism was first described by Albert Neisser in 1879 in urethral exudate. The Thayer-Martin selective medium was developed in the 1960s. NAATs for gonorrhea were introduced in the 1990s and became the recommended diagnostic method by the CDC in the early 2000s. Evolution of antimicrobial resistance in \textit{N. gonorrhoeae} has been relentless: sulfonamide resistance in the 1940s, penicillin resistance (chromosomal and plasmid-mediated penicillinase-producing \textit{N. gonorrhoeae}/PPNG) in the 1970s, tetracycline resistance in the 1980s, fluoroquinolone resistance in the 1990s--2000s (leading the CDC to remove fluoroquinolones as recommended treatment in 2007), and the emergence of decreased susceptibility and resistance to oral cephalosporins (cefixime) in the 2010s. As of 2020, intramuscular ceftriaxone (500 mg for weight <150 kg; 1 g for $\geq$150 kg) is the sole CDC-recommended therapy. The global emergence of extensively drug-resistant (XDR) \textit{N. gonorrhoeae} (ceftriaxone-resistant, azithromycin-resistant) is a WHO-critical priority. The DISPATCH study (2010s) and others demonstrated the efficacy of ceftriaxone monotherapy.

\section*{Why This Test Is Ordered}
\begin{itemize}
  \item Screening for urogenital gonorrhea in sexually active women under 25 years, women 25 and older with risk factors (new or multiple sexual partners, a partner with STI, commercial sex work, inconsistent condom use), and all pregnant women at the first prenatal visit with retesting in the third trimester if at ongoing risk
  \item Screening for urogenital gonorrhea in MSM at least annually (urethral, rectal, and pharyngeal sites based on reported sexual practices); more frequent screening (every 3--6 months) for MSM at high risk (multiple partners, methamphetamine use, partner with STI)
  \item Diagnosis of gonococcal urethritis in men presenting with dysuria, urethral discharge, or urethral pruritus
  \item Evaluation of women with vaginal discharge, intermenstrual bleeding, dyspareunia, or pelvic pain consistent with cervicitis or PID
  \item Evaluation of pharyngeal symptoms (sore throat, pharyngeal erythema) or anorectal symptoms (discharge, tenesmus, pruritus, bleeding) in individuals reporting receptive oral or anal sex without barrier protection
  \item Testing of sexual partners of individuals diagnosed with gonorrhea (expedited partner therapy [EPT] where legally permissible: dispensing cefixime 800 mg PO as a single dose for the partner without medical evaluation)
  \item Diagnosis of disseminated gonococcal infection (DGI) in patients with the classic triad of tenosynovitis, dermatitis, and asymmetric polyarthralgia, or culture-negative septic arthritis in a sexually active individual
  \item Evaluation of neonatal purulent conjunctivitis (ophthalmia neonatorum) by Gram stain, culture, and NAAT of conjunctival exudate
\end{itemize}

\section*{Is This a Primary or Secondary Test}
NAAT is the primary screening and diagnostic test for gonorrhea at all anatomic sites. Culture is a primary test when antimicrobial susceptibility testing is required (treatment failure, suspected resistance) and for pharyngeal and rectal specimens if NAAT is not FDA-cleared for those specific specimen types by the laboratory. Gram stain of male urethral discharge is a primary rapid diagnostic in symptomatic men.

\section*{How This Test Is Performed}
First-catch urine (FCU): The patient should not have urinated for at least 1 hour. Collect the first 10--20 mL of the urine stream (not midstream) in a sterile cup. Vaginal swab (self-collected or clinician-collected): A flocked swab is inserted 2--3 inches into the vagina, rotated for 10--30 seconds, and withdrawn. Vaginal swabs are equivalent to endocervical swabs for NAAT and are preferred for screening because they avoid the need for speculum examination. Endocervical swab: A swab is inserted into the endocervical canal and rotated for 10--30 seconds after removing excess mucus. This requires a speculum examination. Pharyngeal swab: Swab the tonsillar pillars and posterior pharynx; avoid touching the tongue or uvula. Rectal swab: Blindly insert a swab 2--3 cm into the anal canal and rotate; avoid fecal contamination. Conjunctival swab: Swab the purulent exudate and everted eyelid. For culture, specimens should be plated directly onto pre-warmed MTM or NYC agar at the bedside (preferred) or placed immediately into Amies transport medium with charcoal and transported to the laboratory within 24 hours. NAAT specimens are placed in the manufacturer-specific transport medium.

\section*{What Preparation Is Needed}
Patients should not urinate for at least 1 hour before urine collection (to allow adequate accumulation of exfoliated epithelial cells containing organisms). Women should not douche, use vaginal creams, or insert suppositories for 24 hours before vaginal or endocervical specimen collection. Patients should not take antibiotics for at least 2--3 weeks before test of cure (TOC) if performed. Recent use of mouthwash or gargling may transiently reduce pharyngeal gonococcal load and should be avoided for 1 hour before pharyngeal specimen collection.

\section*{Contraindications and Precautions}
NAATs may remain positive for up to 3 weeks after successful treatment due to detection of nonviable nucleic acid. Test of cure (TOC) should be performed no sooner than 7 days after treatment completion, and preferably 2--3 weeks after treatment to avoid false positives. All patients diagnosed with gonorrhea should be treated presumptively for chlamydia co-infection (doxycycline 100 mg PO BID x 7 days) unless chlamydia is excluded by NAAT, given a 10--30\% co-infection rate. Culture has reduced sensitivity compared to NAAT, particularly for asymptomatic and pharyngeal infections. The absence of culture growth in a NAAT-positive specimen does not rule out infection; it may represent low organism burden or loss of viability during transport. Gram stain of asymptomatic male urethral discharge has low sensitivity and should not be used for screening. Pharyngeal gonorrhea is typically asymptomatic (present in 5--15\% of MSM) and serves as a reservoir for transmission and breeding ground for resistance due to the anatomical proximity to commensal \textit{Neisseria} species from which \textit{N. gonorrhoeae} can acquire resistance genes via horizontal gene transfer.

\section*{Risks}
Specimen collection from urethra, cervix, pharynx, or rectum is minimally invasive and carries negligible risk. Urine collection is noninvasive. The primary clinical risk of gonorrhea is, in women, ascending infection to PID with tubal scarring causing infertility (risk increases with each episode) and ectopic pregnancy. In both sexes, DGI and disseminated infection occur in 0.5--3\% of untreated gonorrhea. Untreated gonorrhea also increases HIV acquisition and transmission risk.

\section*{What Happens After the Test}
NAAT results available in 1--4 hours (rapid platforms like GeneXpert: 90 minutes). Culture: preliminary growth at 24--48 hours; final identification and susceptibility at 48--72 hours. Patients diagnosed with gonorrhea should receive ceftriaxone 500 mg IM (1 g IM if weight $\geq$150 kg) as a single dose. If chlamydial co-infection has not been excluded, doxycycline 100 mg PO BID x 7 days should be co-administered. Sexual partners from the preceding 60 days should be evaluated, tested, and treated presumptively. EPT is recommended where legally permissible. Test of cure is NOT routinely recommended for uncomplicated urogenital or pharyngeal gonorrhea treated with the recommended regimen. However, TOC is indicated for pharyngeal gonorrhea (test at 7--14 days), PID, DGI, in pregnancy, or when non-ceftriaxone regimens are used. All cases must be reported to public health.

\section*{Understanding Results}

\subsection*{Below Normal}
\begin{itemize}
  \item \textbf{NAAT negative (not detected):} No \textit{N. gonorrhoeae} nucleic acid detected at the tested site. Essentially excludes gonorrhea at that site with high sensitivity. However, failure to test all exposed anatomic sites can miss infection (e.g., negative urethral test does not exclude pharyngeal or rectal infection)
  \item \textbf{Culture negative:} No \textit{N. gonorrhoeae} isolated. May be a false negative if specimen transport was delayed (gonococci are labile), the patient received recent antibiotics, or the organism burden is low (asymptomatic infection, pharyngeal infection)
  \item \textbf{Gram stain negative (male urethral):} In a symptomatic male with discharge, a negative Gram stain makes gonorrhea unlikely but should be confirmed by NAAT. For extramale urethral specimens, a negative Gram stain does not exclude infection
\end{itemize}

\subsection*{Above Normal}
\begin{itemize}
  \item \textbf{NAAT positive for \textit{N. gonorrhoeae}:} Confirms gonococcal infection at the tested site. All patients with urogenital gonorrhea should be treated with ceftriaxone-based therapy plus doxycycline (for chlamydia coverage). Pharyngeal gonorrhea may require test of cure given lower treatment efficacy at this site and concern for emerging resistance. Rectal and pharyngeal infections may be asymptomatic and serve as ongoing reservoirs for transmission
  \item \textbf{Culture positive for \textit{N. gonorrhoeae} with antimicrobial susceptibility:} Definitive diagnosis. AST results (ceftriaxone MIC, cefixime MIC, azithromycin MIC, tetracycline MIC, ciprofloxacin MIC) are reported. Ceftriaxone MICs $\leq$0.125 mcg/mL are considered susceptible using the EUCAST breakpoint (no CLSI breakpoint exists for ceftriaxone but CLSI recommends reporting results). Elevated ceftriaxone MICs (0.125--0.25 mcg/mL) have been associated with treatment failure at some anatomic sites. The CDC is monitoring ceftriaxone susceptibility through GISP (Gonococcal Isolate Surveillance Project)
  \item \textbf{Gram stain positive (gram-negative intracellular diplococci in PMNs):} In symptomatic males, highly sensitive and specific for gonorrhea. In a female, gram-negative diplococci on an endocervical Gram stain have moderate PPV (~50\%) and must be confirmed by culture or NAAT
  \item \textbf{Positive from normally sterile site (blood, synovial fluid, CSF):} Diagnostic of disseminated gonococcal infection (DGI). DGI blood or synovial fluid cultures are positive in only 10--30\% of cases; mucosal site testing for \textit{N. gonorrhoeae} (urogenital, pharyngeal, rectal) is often positive (50--80\%) even when sterile sites are culture-negative. A clinical diagnosis of DGI should be made in a patient with compatible symptoms and confirmed mucosal \textit{N. gonorrhoeae} infection
\end{itemize}

\section*{Normal Results}
\textbf{\textit{N. gonorrhoeae} not detected} (NAAT) at all tested sites. Culture: \textbf{No \textit{N. gonorrhoeae} isolated}. Urethral Gram stain: \textbf{No gram-negative intracellular diplococci}.

\section*{Follow-up Tests}
\begin{itemize}
  \item \textbf{Gonorrhea culture with antimicrobial susceptibility testing (AST):} When treatment failure is suspected (persistent symptoms or positive NAAT after treatment), when the patient acquired infection in a region with known high-level resistance, or for pharyngeal gonorrhea to document ceftriaxone susceptibility. Culture must be requested explicitly (not all laboratories perform gonococcal AST)
  \item \textbf{Test of cure (NAAT):} For pharyngeal gonorrhea at 7--14 days after treatment; for all patients treated with a non-ceftriaxone regimen; for pregnant women; for patients with PID or DGI. NAAT can be performed as early as 7 days but ideally 14 days post-treatment
  \item \item \textit{\textbf{Chlamydia trachomatis}} \textbf{NAAT:} Co-testing for chlamydia from the same specimens (most NAATs are combination CT/NG) at all exposed sites; co-infection rates are 10--30\%
  \item \textbf{HIV testing, syphilis serology, hepatitis B and C screening:} Comprehensive STI evaluation is recommended given the shared risk factors
  \item \textbf{Pelvic ultrasound:} For women diagnosed with gonorrhea who have pelvic pain, adnexal tenderness, or cervical motion tenderness to evaluate for tubo-ovarian abscess, a complication of PID
\end{itemize}

\section*{References}
\begin{enumerate}
  \item Workowski KA, Bachmann LH, Chan PA, et al. Sexually transmitted infections treatment guidelines, 2021. MMWR Recomm Rep. 2021;70(4):1-187.
  \item Unemo M, Lahra MM, Cole M, et al. World Health Organization Global Gonococcal Antimicrobial Surveillance Program (WHO GASP). Sex Transm Infect. 2023; epub.
  \item CDC. Update to CDC's treatment guidelines for gonococcal infection, 2020. MMWR Morb Mortal Wkly Rep. 2020;69(50):1911-1916.
  \item Golden MR, Whittington WL, Handsfield HH, et al. Effect of expedited treatment of sex partners on recurrent or persistent gonorrhea or chlamydial infection. N Engl J Med. 2005;352(7):676-685.
  \item Kirkcaldy RD, Harvey A, Papp JR, et al. Neisseria gonorrhoeae antimicrobial susceptibility surveillance --- The Gonococcal Isolate Surveillance Project, 27 Sites, United States, 2014. MMWR Surveill Summ. 2016;65(7):1-19.
\end{enumerate}

\clearpage
